본 탐색은 원시 YOLO의 코너 0--7에 pose-coupled 변위를 더하고 center8·박스·점수·누락 마스크를 유지한다. NoOp는 9개 원시점의 정확한 복사다. 후보 생성의 자세 $T_k$가 아니라 실제 전체 prediction-only W/D+PnP 출력 $F(q_k)$를 평가한다. 공동 bank의 최선 후보 oracle은 그 bank 안의 단일 action 선택 J의 상한이며, 각 코너가 독립 선택하는 I는 bank 밖의 조합도 만들 수 있다. 두 readout을 개발 성능으로 재선택하지 않는다.

GEO의 전체 유효 action 집합을 먼저 고정한 뒤 PERM은 non-NoOp 행의 코너별 대응만 고정 permutation으로 바꾸는 유한 대조다. Disentangling by Factorising (2018/ICML)의 차원별 permutation 원리\cite{factorvae2018}를 참고하되 해당 논문의 FactorVAE가 팔레트 개선을 입증하거나 PERM이 완전한 통계 독립을 보장한다고 주장하지 않는다. 같은 점별 후보 좌표·특징·metadata 주변 multiset과 마스크, 노출량을 검증한다.

후보 생성은 초기 prediction-only W/D 가설별 카메라 축 회전·이동의 12개 부호 축과 3개 반경, 64개 여섯 축 부호 조합을 사용한다. 수치 투영 민감도로 회전(rad)과 이동(m)의 크기를 정규화하고 실제 비선형 투영의 root search로 이동 상한을 맞춘다. 최대 후보는 NoOp를 포함해 201개다. $R'=\operatorname{Exp}(\delta r)R$, $t'=t+\delta t$로 물체 중심에 회전을 적용하며 $\operatorname{Exp}(\delta r)t$로 이동을 회전시키지 않는다. 이동은 영상 대각선의 1\% 이내이고 각 코너를 따로 clipping하지 않는다. 등록 치수와 카메라·왜곡 계약은 기존 실제 prediction-only 경로와 같으며, 후보 생성 및 배포 선택에 자세 참조를 넣지 않는다. 합성 치수는 TRAIN 렌더링 side table의 실제 특권 기하이고 실사는 알려진 객체 종류의 등록 치수여서 같은 정보 취득 경로라고 가정하지 않는다.

코너 $j$의 후보 $k$ 점수 $\ell_{jk}$와 공통 추론 지원 집합 $V$에서 I는 각 $j$의 hard argmax를, J는 $\arg\max_k |V|^{-1}\sum_{j\in V}\ell_{jk}$ 한 개를 사용한다. 동률은 NoOp를 우선하며 정답 가시성이 아닌 고정된 추론 마스크로 점수를 모은다. 동결 실험은 기존 N3의 세 seed 특징·점수기를 그대로 재사용하여 I/J를 모두 보고한다. 새 후보로의 분포 변화가 있으므로 동결 점수기 실패를 특징에 정보가 없다는 증명으로 해석하지 않는다. 추가 학습은 같은 초기 가중치·예시 순서·구조·파라미터 수(20,259개)의 GEO/PERM 쌍을 각 세 seed로 시작한다. synthetic TRAIN 55,980행 중 기존 매칭·목표 지원 규칙을 만족한 55,915행을 사용하고 65행은 제외한다. batch 16, AdamW, seed당 6,000 update·96,000 노출과 마지막 체크포인트를 고정하며 실사 학습·추가 cap/seed/epoch 선택은 없다. 목표는 원시 예측에서 고정한 허용 GT 대칭 대응의 코너 제곱 오차로 만든 후보 분포이고, 공통 지원의 점수 평균과 cross entropy를 사용한다. source selection 1,031장은 진입 oracle 진단, 이미 소비한 합성 1,985장은 개발 진단, 실사 319장/13세션은 재사용 DEV로 서로 분리한다.

source selection 1,031장의 GEO bank에서 실제 전체 자세 함수 $F$의 대칭 8코너 거리(ADDsym) oracle 여지가 수치오차 $10^{-7}$m보다 큰 프레임은 954장이었고, 여지 중앙값은 0.011915m였다. 원시와 oracle 모두 1,030/1,031장에서 자세를 산출했다. 이 상한은 참조로 고른 개발 진단이며 배포 결과가 아니다. 이에 따라 GEO/PERM 각각 세 seed를 원래 무작위 초기화부터 6,000 update씩, 총 36,000 update와 576,000 노출로 학습했다. smoke 4 update는 별도로 폐기했다.

\begin{table*}[t]\centering\small\caption{실제 DEV의 추가 탐색 실험. 세 seed의 원결과; 코너 중앙값/90백분위는 매칭 관측 코너의 조건부 통계이며 PCK10은 전체 참조 코너를 분모로 한 0--1 분율이다. 코너 오차는 311장/2,445개 관측 코너, PCK는 2,499개 참조 코너, 위치·회전은 319/319장 자세 산출을 사용한다. 자세는 같은 2D 레이블에서 재구성한 참조를 사용한다.}\label{sup:joint_action_results}\begin{tabular}{llrrrrr}\toprule 방법&seed&med(px)&P90(px)&PCK10&T(cm)&R(deg)\\\midrule

GEO--J & 1 & 6.429 & 44.777 & 0.648 & 7.289 & 2.403\\

GEO--J & 2 & 6.461 & 43.095 & 0.649 & 7.663 & 2.467\\

GEO--J & 3 & 6.491 & 43.986 & 0.642 & 7.885 & 2.492\\

PERM--J & 1 & 6.627 & 44.105 & 0.637 & 8.227 & 2.550\\

PERM--J & 2 & 6.766 & 43.590 & 0.638 & 8.324 & 2.558\\

PERM--J & 3 & 6.564 & 43.675 & 0.643 & 7.897 & 2.358\\

\bottomrule\end{tabular}\end{table*}

GEO--J minus N3의 frame별 seed 평균 정규화 코너 거리 차이는 0.00078570, 95\% session bootstrap 구간은 [0.00045428, 0.00125593]이고 고정 손상/보존 판정은 오차 악화였다.

GEO--J minus PoseFix의 frame별 seed 평균 정규화 코너 거리 차이는 0.00090093, 95\% session bootstrap 구간은 [0.00050038, 0.00141278]이고 고정 손상/보존 판정은 오차 악화였다.

GEO--J minus RAW의 frame별 seed 평균 정규화 코너 거리 차이는 -0.00031911, 95\% session bootstrap 구간은 [-0.00065231, -0.00005401]이고 고정 손상/보존 판정은 개선·보존 조건 미충족이었다.

GEO--J minus FIT--PERM--J의 frame별 seed 평균 정규화 코너 거리 차이는 -0.00019133, 95\% session bootstrap 구간은 [-0.00039871, -0.00002461]이고 고정 손상/보존 판정은 개선·보존 조건 미충족이었다.

실제 319장/13세션은 반복 사용한 개발 자료다. 모든 319장에 검출과 유효한 8코너가 존재하지만 GT instance 매칭은 311장이다. 실패와 매칭을 분리했고, 원결과/seed별 손상/같은 bank oracle gap 및 synthetic 개발 진단은 실행 보고서에 남겼다. GEO/PERM의 결합 대조와 동결 scorer의 I/J 대조는 기전의 제한된 범위만 다루며 N3/PoseFix 대비 방법 전체의 효과와 구분한다. 새로운 독립 참조의 clean/실물 가림 pair가 없으므로 독립적인 센서 6D 정확도 확증은 완료하지 못했다. 배포 runtime은 동일 입력/장치/전체 경계의 별도 영수증과 통합한다. 기존 N3 headline은 변경하지 않았다.


\subsection{대조별 효과와 손상·보존}

GEO--J의 세 seed 코너 중앙값은 모두 기존 N3(세 seed 통계 평균 5.778px) 및 PoseFix(5.561px)보다 높았다. 이 중앙값 비교와 $E_{\rm sym}$의 frame별 paired 평균 차이는 서로 다른 통계량이다. 표~\ref{sup:joint_action_damage}는 정확한 정준 GT 코너를 연결한 심한 손상을 보존한다. RAW/PERM 대비 정규화 오차 구간이 음수여도 기존 좋은 코너를 보존하는 고정 조건을 통과하지 못했다. 관측된 이득을 단일 양의 연구 결론으로 합치지 않는다.

\begin{table*}[t]\centering\small\caption{GEO--J minus 대조의 실제 DEV 손상과 자세 변화. 손상은 정준 GT의 같은 참조 코너(전체 2,499개; 매칭 관측 2,445개)에서 대조 오차 $<5$px가 GEO 오차 $>10$px가 된 개수이며 seed 1/2/3을 각각 표시한다. 이동·회전은 세 seed 모두 공동 성공한 319장의 paired 평균 차이이며 표의 위치·회전 중앙값 차이와 다르다. 구간은 고정 13세션/10,000회 bootstrap의 탐색적 95\% 구간이다.}\label{sup:joint_action_damage}\begin{tabular}{lrrr}\toprule 대조 & 손상 코너(seed 1/2/3) & 이동 차이(cm), 95\% 구간 & 회전 차이(deg), 95\% 구간 \\\midrule

RAW & 6/6/5 & -0.324 [-0.859, 0.027] & -0.060 [-0.236, 0.080]\\

N3 & 17/15/19 & -0.131 [-3.232, 1.441] & 2.909 [0.820, 4.780]\\

PoseFix & 28/26/26 & -0.469 [-3.224, 1.058] & 3.293 [0.994, 5.441]\\

PERM--J & 5/15/22 & -0.598 [-1.379, -0.148] & -0.107 [-0.963, 0.418]\\

\bottomrule\end{tabular}\end{table*}

GEO--J minus N3와 PoseFix의 paired 회전 평균 차이 구간은 각각 [0.820, 4.780]도와 [0.994, 5.441]도로 악화를 보였다. 두 이동 평균 차이 구간은 0을 포함하며 독립 물리 정확도 확증으로 해석하지 않는다. GEO/PERM 모두 세 seed에서 최종 자세를 319/319장 산출했고 2D 매칭은 311장이었다. 보존된 박스·점수·마스크와 자세 산출률은 원래 정확한 코너·자세가 보존되었다는 뜻이 아니다.

\begin{table*}[t]\centering\small\caption{실제 DEV의 고정 J 선택과 같은 bank oracle 차이. NoOp/최종 W--D 가지 변경은 319장 중 개수이고, oracle gap은 같은 319장 공동 성공의 $\mathrm{ADD}_{\rm sym}(J)-\mathrm{ADD}_{\rm sym}(oracle)$ 중앙값(m)이다. 생성 가설과 최종 자세 함수 $F$의 가지는 구별한다. PERM은 일관된 생성 자세가 없어서 가지 변경의 비교값을 NA로 둔다.}\label{sup:joint_action_selection}\begin{tabular}{llrrr}\toprule 방법 & seed & NoOp & 최종 W--D 변경 & oracle gap med(m) \\\midrule

GEO--J & 1 & 82 & 117 & 0.029069\\

GEO--J & 2 & 100 & 100 & 0.027060\\

GEO--J & 3 & 129 & 103 & 0.030040\\

PERM--J & 1 & 140 & NA & 0.027392\\

PERM--J & 2 & 121 & NA & 0.029288\\

PERM--J & 3 & 109 & NA & 0.026669\\

\bottomrule\end{tabular}\end{table*}

source oracle에서 PERM도 비영 여지 1,003장과 여지 중앙값 0.013282m로 GEO보다 큰 후보 상한을 가졌다. 따라서 GEO oracle의 존재만으로 결합의 이득을 증명하지 않는다. 실사에서 GEO--J minus PERM--J의 음의 정규화 오차 구간은 고정 손상 조건을 통과하지 못했고, 이미 소비한 합성 1,985장 진단에서는 같은 차이가 0.00010170이며 95\% frame bootstrap 구간 [0.00007133, 0.00013190]로 GEO가 악화되었다. 합성 정규화 코너 오차의 paired 구간은 같은 참조 지원이 있는 1,978/1,985장이다. 7장은 참조 코너가 없어 non-evaluable로 명시하며, 전체 자세 산출은 1,985/1,985장이다. 합성 주 분석은 frame 재표집을 실행했으며 선택적인 scenario 보조 재표집은 실행하지 않았다.

동결 GEO의 J minus I 정규화 오차 차이는 0.00006435, 실사 95\% session 구간 [-0.00001740, 0.00018203]로 0을 포함했다. PERM의 해당 차이도 0.00006863, 구간 [-0.00000907, 0.00018156]로 0을 포함했다. GEO의 I NoOp(8코너 모두 index0)는 seed별147/141/167장, J는294/308/307장이었다. 각 I는 bank 밖의 조합도 만들 수 있으므로 J의 oracle gap을 I의 상한 위반으로 읽지 않는다.

최신 사람 등급 153/92/74장, 기존 거리 tag 155/103/59장과 unknown2장, 2,499개 사람 코너 상태의 보조 재집계는 원자료와 같은 GT 대응으로 별도 결과에 보존하였다. 등급 기준은 미확인이고 거리 tag는 독립 계측이 아니며 any-visibility 프레임 집단은 겹친다. 이를 물리 가림 정도의 인과 효과나 독립 거리별 정확도로 확대하지 않는다. 모든 구간은 반복 DEV의 탐색적 분석이고 다중비교 보정을 적용하지 않았다. 후보·동결 특징·목표·예산의 이 제한 실험은 용량 부족이나 상관된 코너 오류가 지배 원인이라는 진단을 확정하지 않는다.


\subsection{추가 후보 방법의 실제 배포 경계 비용}

표~\ref{sup:joint_action_runtime}는 위 정확도 평가와 같은 seed1 가중치·hard J 경로를 새로운 조용한 계측 패널에서 비교한다. 본문의 Base/N3/PoseFix 패널과 같은 장치·경계지만 별도의 실행이므로 Base/N3도 이 패널의 원값을 따로 보인다. 모든 A 호출은 RAM의 native BGR 영상에서 시작하여 전처리·검출·공유 특징, 초기 prediction-only W/D 가설·민감도와 실제 투영 root search 후보 생성, 새로운 RGB 특징 patch 점수·hard J, 선택된 점의 실제 최종 W/D+PnP $F(q)$를 포함한다. 파일 decoding/loading과 oracle의 모든 후보 최종 $F$ 탐색은 포함하지 않는다.

\begin{table}[t]\centering\small\caption{같은 RTX3080/pallet-yolo26 환경에서 새 A 계측 패널의 전체 벽시계 시간(ms). 13세션/26프레임, arm별 준비20회·프레임당5반복, batch1·사전 고정 seed1이다. 후보 arm은 각201후보를 실제 생성·채점하며 메모리는 모든 모델이 상주한 process여서 방법별 독립 메모리 순위가 아니다.}\label{sup:joint_action_runtime}\begin{tabular}{lrrr}\toprule 방법 & 중앙값(ms) & P90(ms) & 자세 산출 \\\midrule

Base & 11.080 & 13.428 & 130/130\\

N3 & 14.433 & 16.636 & 130/130\\

동결 GEO--J & 56.663 & 59.235 & 130/130\\

학습 GEO--J & 56.541 & 59.174 & 130/130\\

학습 PERM--J & 56.667 & 59.837 & 130/130\\

\bottomrule\end{tabular}\end{table}

측정650회와 준비100회 총750호출의 원출력 검산에서 모든 arm의 최대 좌표 차이는 0px였고, J action·최종 W/D 가지·자세 산출도 정확히 일치했다. 기존 YOLO/N3의 cuDNN TF32 on·matmul TF32 off 수치 계약을 유지했으며 추가 학습은 없다. 먼저 CPU 집계와 겹쳤던750호출은 비용 비교에서 전부 폐기하고 조용한 재측정만 표에 사용했다. 이 실제 비용에서는 새 후보 방법이 N3보다 느리고 앞의 실제 DEV 정확도·손상도 실용 기준 대비 이득을 지지하지 않으므로 배포 개선으로 주장하지 않는다. 후보 oracle의 오프라인 호출량과 배포 비용을 합산하지 않는다.

